% Source: https://artofproblemsolving.com/community/c3h2139765_rationalization
% Stable identity: aops-topic-2139765-q01
Group the denominator as $\left(1+\sqrt{15}\right)+\sqrt{23}$ and multiply numerator and denominator by $\left(1+\sqrt{15}\right)-\sqrt{23}$:
\[\left(1+\sqrt{15}\right)^2-23=16+2\sqrt{15}-23=2\sqrt{15}-7 .\]
Next multiply by $2\sqrt{15}+7$: $\left(2\sqrt{15}\right)^2-7^2=60-49=11$. Since $2\sqrt{12}=4\sqrt3$,
\[\dfrac{2\sqrt{12}}{1+\sqrt{15}+\sqrt{23}}=\dfrac{4\sqrt3\left(1+\sqrt{15}-\sqrt{23}\right)\left(2\sqrt{15}+7\right)}{11}.\]
Expanding,
\[\left(1+\sqrt{15}-\sqrt{23}\right)\left(2\sqrt{15}+7\right)=2\sqrt{15}+7+30+7\sqrt{15}-2\sqrt{345}-7\sqrt{23}=37+9\sqrt{15}-7\sqrt{23}-2\sqrt{345},\]
and multiplying by $4\sqrt3$, with $\sqrt{45}=3\sqrt5$ and $\sqrt{1035}=3\sqrt{115}$,
\[4\sqrt3\left(37+9\sqrt{15}-7\sqrt{23}-2\sqrt{345}\right)=148\sqrt3+108\sqrt5-28\sqrt{69}-24\sqrt{115}.\]
So
\[\dfrac{2\sqrt{12}}{1+\sqrt{15}+\sqrt{23}}=\boxed{\dfrac{148\sqrt3+108\sqrt5-28\sqrt{69}-24\sqrt{115}}{11}}\approx0.7165 .\]
